% !TEX root = main.tex
% List of Symbols (included via \symbols{% !TEX root = main.tex
% List of Symbols (included via \symbols{% !TEX root = main.tex
% List of Symbols (included via \symbols{% !TEX root = main.tex
% List of Symbols (included via \symbols{\input{listofsymbols}} in main.tex)
%
% A two-column table: the symbol, then its description. Start each
% description with a capital letter and write it in sentence case (GSO).
% Organize the list by category or function, alphabetically by name, or by
% frequency of use (GSO). Alphabetical order follows the characters
% themselves, so * comes after &. The sample is organized by category.
% The table continues across pages automatically. To widen the symbol
% column, give a width: \begin{symboltable}[1.5in]

\begin{symboltable}
% Number systems
\sym{$\ZZ$}{The ring of integers}
\sym{$\RR$}{The field of real numbers}
\sym{$\CC$}{The field of complex numbers}
% Spaces and groups
\sym{$H^1(\RR)$}{The Sobolev space of functions over $\RR$ that are square-integrable and whose weak first derivative is square-integrable}
\sym{$\mathrm{GL}_2(\RR)$}{The general linear group of real invertible $2\times 2$ matrices}
% Physical quantities
\sym{$c$}{Speed of light in vacuum ($c \approx 3 \times 10^8\text{ m/s}$)}
\sym{$E$}{Energy of a physical system}
\sym{$m$}{Mass of a physical body}
% Other notation
\sym{$\%$}{Percentage sign}
\end{symboltable}
} in main.tex)
%
% A two-column table: the symbol, then its description. Start each
% description with a capital letter and write it in sentence case (GSO).
% Organize the list by category or function, alphabetically by name, or by
% frequency of use (GSO). Alphabetical order follows the characters
% themselves, so * comes after &. The sample is organized by category.
% The table continues across pages automatically. To widen the symbol
% column, give a width: \begin{symboltable}[1.5in]

\begin{symboltable}
% Number systems
\sym{$\ZZ$}{The ring of integers}
\sym{$\RR$}{The field of real numbers}
\sym{$\CC$}{The field of complex numbers}
% Spaces and groups
\sym{$H^1(\RR)$}{The Sobolev space of functions over $\RR$ that are square-integrable and whose weak first derivative is square-integrable}
\sym{$\mathrm{GL}_2(\RR)$}{The general linear group of real invertible $2\times 2$ matrices}
% Physical quantities
\sym{$c$}{Speed of light in vacuum ($c \approx 3 \times 10^8\text{ m/s}$)}
\sym{$E$}{Energy of a physical system}
\sym{$m$}{Mass of a physical body}
% Other notation
\sym{$\%$}{Percentage sign}
\end{symboltable}
} in main.tex)
%
% A two-column table: the symbol, then its description. Start each
% description with a capital letter and write it in sentence case (GSO).
% Organize the list by category or function, alphabetically by name, or by
% frequency of use (GSO). Alphabetical order follows the characters
% themselves, so * comes after &. The sample is organized by category.
% The table continues across pages automatically. To widen the symbol
% column, give a width: \begin{symboltable}[1.5in]

\begin{symboltable}
% Number systems
\sym{$\ZZ$}{The ring of integers}
\sym{$\RR$}{The field of real numbers}
\sym{$\CC$}{The field of complex numbers}
% Spaces and groups
\sym{$H^1(\RR)$}{The Sobolev space of functions over $\RR$ that are square-integrable and whose weak first derivative is square-integrable}
\sym{$\mathrm{GL}_2(\RR)$}{The general linear group of real invertible $2\times 2$ matrices}
% Physical quantities
\sym{$c$}{Speed of light in vacuum ($c \approx 3 \times 10^8\text{ m/s}$)}
\sym{$E$}{Energy of a physical system}
\sym{$m$}{Mass of a physical body}
% Other notation
\sym{$\%$}{Percentage sign}
\end{symboltable}
} in main.tex)
%
% A two-column table: the symbol, then its description. Start each
% description with a capital letter and write it in sentence case (GSO).
% Organize the list by category or function, alphabetically by name, or by
% frequency of use (GSO). Alphabetical order follows the characters
% themselves, so * comes after &. The sample is organized by category.
% The table continues across pages automatically. To widen the symbol
% column, give a width: \begin{symboltable}[1.5in]

\begin{symboltable}
% Number systems
\sym{$\ZZ$}{The ring of integers}
\sym{$\RR$}{The field of real numbers}
\sym{$\CC$}{The field of complex numbers}
% Spaces and groups
\sym{$H^1(\RR)$}{The Sobolev space of functions over $\RR$ that are square-integrable and whose weak first derivative is square-integrable}
\sym{$\mathrm{GL}_2(\RR)$}{The general linear group of real invertible $2\times 2$ matrices}
% Physical quantities
\sym{$c$}{Speed of light in vacuum ($c \approx 3 \times 10^8\text{ m/s}$)}
\sym{$E$}{Energy of a physical system}
\sym{$m$}{Mass of a physical body}
% Other notation
\sym{$\%$}{Percentage sign}
\end{symboltable}
